\section{Nonseparable reflection jets: a two-generator theorem}
\label{ns:sec:jets}

The manuscript proves a multivariable product law when the active deformation
parameters are separate powers of coordinate variables. Its warning that a
general multivariable substitution requires the full Koszul complex is
essential. We now give a different sufficient condition: the ideal generated
by the active parameters may have arbitrary mixed terms, provided it needs at
most two generators. We also exhibit an explicit three-generator obstruction.
The results concern the specified formal distribution presentation. They do
not assert independence of evaluated polylogarithm or zeta values.

\subsection{The distribution presentation and the input theorem}

For an integer $q>2$, put $X_q=(q^{-1}\mathbb Z)/\mathbb Z$. Given weights
$a_p$ in a commutative coefficient ring $S$, one for each prime $p\mid q$,
define
\begin{equation}\label{ns:eq:distribution}
 \mathcal D_q(S,a)=
 \left(\bigoplus_{x\in X_q}S e_x\right)
 \Big/
 \left\langle
  \sum_{\substack{y\in X_q\\py=x}}e_y-a_pe_x:
  p\mid q,\ x\in X_{q/p}
 \right\rangle.
\end{equation}
Reflection acts by $ce_x=e_{-x}$. Write
\[
 n=\frac{\varphi(q)}2,\qquad
 r=\#\{p>2:p\mid q\}+\mathbf1_{4\mid q}.
\]
The active sequence is
\begin{equation}\label{ns:eq:active}
 \boldsymbol\delta=
 (a_p-1:p>2,\ p\mid q),
 \quad\text{with }a_2(a_2-1)\text{ appended if }4\mid q.
\end{equation}
There is no active entry at $2$ when $v_2(q)=1$.

We use the following established results from the frozen manuscript \cite{repoJets}.
The all-ring basis theorem makes $\mathcal D_q(S,a)$ free of rank
$\varphi(q)$. In characteristic two, with $N=c-1$, its reflection homology is
\begin{equation}\label{ns:eq:input-koszul}
 \ker N/\operatorname{im}N
 \simeq\bigoplus_{j=0}^r H_j(K_S(\boldsymbol\delta)).
\end{equation}
For a finite free $\mathbb Z$-algebra $B$, put $\overline B=B/2B$ and
$C_\pm=\mathcal D_q(B,a)/(c\mp1)\mathcal D_q(B,a)$. Then
\begin{align}
 \operatorname{Tor}_{\mathbb Z}C_+
 &\simeq\bigoplus_{j\ {
m odd}}
 H_j(K_{\overline B}(\overline{\boldsymbol\delta})),
 \label{ns:eq:input-plus}\\
 \operatorname{Tor}_{\mathbb Z}C_-
 &\simeq\bigoplus_{j\ {
m even}}
 H_j(K_{\overline B}(\overline{\boldsymbol\delta})).
 \label{ns:eq:input-minus}
\end{align}
Every integer-torsion class is killed by two, and both free abelian ranks are
$n\operatorname{rank}_{\mathbb Z}B$. The source labels are
\path{ir:thm:integral} in
\path{08-integral-02-integral.tex},
\path{ir:thm:char2-koszul} in
\path{08-integral-05-modular-jets.tex}, and
\path{irjet:thm:finite-free-koszul} in
\path{08-integral-jets.tex}. Thus the algebraic input and the new step below
are separately identifiable.

\subsection{A colength law without separated parameters}

A finite-dimensional commutative $k$-algebra $S$ is \emph{Frobenius} here if it
admits a $k$-linear functional $\tau:S\to k$ such that $(u,v)\mapsto\tau(uv)$
is a nondegenerate pairing. The rectangular jet algebra
\begin{equation}\label{ns:eq:rectangular}
 S=k[t_1,\ldots,t_d]/(t_1^{L_1},\ldots,t_d^{L_d}),\qquad L_i\ge1,
\end{equation}
is a local Frobenius algebra. Take $\tau$ to be the coefficient of
$t_1^{L_1-1}\cdots t_d^{L_d-1}$: complementary basis monomials pair to one.
Its dimension is $T=\prod_iL_i$.

\begin{theorem}[Two-generator colength law]\label{ns:thm:two-generator}
Let $S$ be a finite-dimensional commutative local Frobenius algebra over a
field $k$. Let $\boldsymbol f=(f_1,\ldots,f_r)$, with $r\ge1$, and suppose
that $I=(f_1,\ldots,f_r)$ can be generated by at most two elements. Put
$\ell=\dim_k(S/I)$. Then
\begin{equation}\label{ns:eq:binomial}
 \dim_k H_j(K_S(\boldsymbol f))=\binom rj\ell,
 \qquad 0\le j\le r.
\end{equation}
This statement holds in every characteristic.
\end{theorem}

\begin{proof}
If $I=S$, some entry is a unit because $S$ is local. Its two-term Koszul
factor is contractible, and both sides of \eqref{ns:eq:binomial} vanish.
If $I=0$, the differential is zero and the assertion is immediate.

Otherwise choose a minimal generating subset among the given entries.
One obtains it by choosing a basis of $I/\mathfrak mI$ from their residue
classes and applying Nakayama's lemma. This subset has at most two elements.
Permute those elements to the beginning of the list. Every remaining entry
is a linear combination of them, so elementary invertible operations change
the list into $(f,g,0,\ldots,0)$, allowing $g=0$ if only one generator is
needed. These operations induce isomorphisms on the free module defining
the Koszul complex and on all its exterior powers.

For $r=1$, the kernel and cokernel of the square linear map
$S\xrightarrow{f}S$ have the same dimension, proving the assertion. For
$r\ge2$, first consider the two-entry complex
\[
 0\longrightarrow S\longrightarrow S^2\longrightarrow S\longrightarrow0.
\]
Its zeroth homology is $S/I$, and its second homology is
$\operatorname{Ann}_S(I)$. The Frobenius pairing identifies
\begin{equation}\label{ns:eq:dual}
 \operatorname{Ann}_S(I)\simeq\operatorname{Hom}_k(S/I,k)
\end{equation}
as $S$-modules. Indeed, the orthogonal complement of $I$ is precisely its
annihilator. To see the nontrivial inclusion, suppose $\tau(zi)=0$ for
all $i\in I$. For every $a\in S$, one then has $\tau((zi)a)=0$, because
$ia\in I$. Nondegeneracy gives $zi=0$.

The two end homology groups therefore have dimension $\ell$. The Euler
characteristic of the complex is $T-2T+T=0$, so its middle homology has
dimension $2\ell$. Tensoring with the zero-differential exterior algebra
on the remaining $r-2$ generators yields
\[
 \dim_kH_j
 =\ell\left[\binom{r-2}j+2\binom{r-2}{j-1}
                  +\binom{r-2}{j-2}\right]
 =\ell\binom rj.
\]
No separability of the entries has been used.
\end{proof}

\begin{corollary}[Nonseparable reflection formula]\label{ns:cor:reflection}
Let $S$ be as in Theorem~\ref{ns:thm:two-generator}, now in characteristic
two, and put $T=\dim_kS$. If the active ideal generated by
\eqref{ns:eq:active} needs at most two generators, then
\begin{equation}\label{ns:eq:reflection}
 \dim_k\mathcal D_q(S,a)/(c-1)\mathcal D_q(S,a)
 =nT+2^{r-1}\ell.
\end{equation}
For
$B=\mathbb Z[t_1,\ldots,t_d]/(t_1^{L_1},\ldots,t_d^{L_d})$, if the reduced
active ideal in $S=B/2B$ has that property, both signs have the abelian-group
decomposition
\begin{equation}\label{ns:eq:smith}
 C_\pm\simeq\mathbb Z^{nT}\oplus
                  (\mathbb Z/2\mathbb Z)^{2^{r-1}\ell}.
\end{equation}
No splitting as $B$-modules is asserted.
\end{corollary}

\begin{proof}
Since $N^2=0$ in characteristic two, an elementary rank calculation gives
\[
 \dim_k\operatorname{coker}N
 =\frac{\dim_k\mathcal D_q(S,a)+
                \dim_k(\ker N/\operatorname{im}N)}2.
\]
The first dimension is $2nT$. Equations \eqref{ns:eq:input-koszul} and
\eqref{ns:eq:binomial} make the second $2^r\ell$, proving
\eqref{ns:eq:reflection}. For the integral assertion, both parity sums of
the binomial coefficients are $2^{r-1}$; apply
\eqref{ns:eq:input-plus}--\eqref{ns:eq:input-minus} and the stated free rank.
\end{proof}

In particular, when $r(q)=2$, this settles the dimension and abelian Smith
problem for \emph{every} active sequence over a rectangular jet algebra.
The number of deformation variables need not equal the active-prime count.
When $r(q)>2$, the same theorem applies whenever the extra active parameters
are redundant generators of the ideal.

\subsection{Mixed monomials and a module obstruction}

Take $S=k[x,y]/(x^L,y^M)$ and
\[
 I=(x^a,x^by^c),\qquad 1\le b<a\le L,\quad1\le c\le M.
\]
A quotient basis consists of the monomials $x^iy^j$ with $i<a$ and with
either $i<b$ or $j<c$. Consequently
\begin{equation}\label{ns:eq:mixed-colength}
 \ell=bM+(a-b)c.
\end{equation}
This includes boundary cases where one displayed monomial is zero in $S$.
For two active parameters, the reflected dimension is therefore
\[
 nLM+2\bigl[bM+(a-b)c\bigr].
\]

For a concrete integral example, let $q=15$,
$B=\mathbb Z[x,y]/(x^3,y^3)$, and take
\[
 a_3=1+x^2,\qquad a_5=1+xy.
\]
Here $T=9$, $r=2$, $n=4$, and $\ell=4$. Thus
\begin{equation}\label{ns:eq:fifteen}
 C_\pm\simeq\mathbb Z^{36}\oplus(\mathbb Z/2\mathbb Z)^8,
 \qquad
 \dim_{\mathbb F_2}\mathcal D_{15}(S,a)/(c-1)=44.
\end{equation}

The module statement of the manuscript's separable product law does not
extend by replacing its quotient with an arbitrary $S/I$. In this example,
put $Q=S/(x^2,xy)$, over $k=\mathbb F_2$. Direct monomial calculation gives
\begin{align*}
 Q&:\quad 1,x,y,y^2,\\
 \operatorname{Ann}_S(I)&=(x^2,xy^2):\quad
             x^2,x^2y,x^2y^2,xy^2.
\end{align*}
The first module is cyclic. The second needs two generators: modulo
$\mathfrak m\operatorname{Ann}_S(I)$, the classes of $x^2$ and $xy^2$ are
independent. Thus \eqref{ns:eq:input-minus} gives
\begin{equation}\label{ns:eq:module-obstruction}
 \operatorname{Tor}_{\mathbb Z}C_-
 \simeq Q\oplus\operatorname{Ann}_S(I),
 \qquad
 \operatorname{Tor}_{\mathbb Z}C_-\not\simeq Q^2.
\end{equation}
The left-hand module requires three generators and the right-hand candidate
requires two. Frobenius duality preserves length; it does not identify an
arbitrary finite module with its dual. This explains why the dimension law
survives although that stronger module formula fails.

For comparison, one may introduce a redundant third active parameter at
$q=105$ by taking
\[
 a_3=1+x^2,\qquad a_5=1+xy,\qquad a_7=1+x^2y+xy^2
\]
over the same ring. The active ideal is unchanged. The homology dimensions
become $(4,12,12,4)$, the integral free rank is $216$, the torsion multiplicity
is $16$, and the binary reflected dimension is $232$.

\subsection{A sharp boundary: three generators can change the answer}

\begin{proposition}[A three-generator obstruction]\label{ns:prop:counterexample}
Let
\[
 S=\mathbb F_2[x,y,z]/(x^2,y^2,z^2),\qquad
 \boldsymbol f=(xy,xz,yz).
\]
Then $\dim_kS=8$, $\dim_kS/(\boldsymbol f)=4$, and
\begin{equation}\label{ns:eq:counter-homology}
 \bigl(\dim H_0,\dim H_1,\dim H_2,\dim H_3\bigr)
 =(4,14,14,4).
\end{equation}
In particular the total homology dimension is $36$, whereas the unsupported
colength extrapolation $2^3\dim_kS/(\boldsymbol f)$ would give $32$.
\end{proposition}

\begin{proof}
The first differential has image spanned by $xy,xz,yz,xyz$, and thus rank
four. With exterior basis vectors denoted by $e_i,e_{ij},e_{123}$, the third
differential has image spanned by
\[
 yz e_{12}+xz e_{13}+xy e_{23},\quad
 xyz e_{12},\quad xyz e_{13},\quad xyz e_{23}.
\]
These vectors are independent, so its rank is also four.

The middle differential has image spanned by
\begin{gather*}
 xz e_1+xy e_2,\qquad yz e_1+xy e_3,\qquad yz e_2+xz e_3,\\
 xyz e_1,\qquad xyz e_2,\qquad xyz e_3.
\end{gather*}
The first three are independent in degree two, and the last three are
independent in degree three. Multiplying any of the first three by $x,y,z$
produces a vector in the span of the final three or zero; higher products
vanish. This list therefore spans the image and has rank six. The chain
dimensions are $8,24,24,8$, so the differential ranks $4,6,4$ give exactly
\eqref{ns:eq:counter-homology}.
\end{proof}

At $q=105$, use the integral algebra
$B=\mathbb Z[x,y,z]/(x^2,y^2,z^2)$ and weights
\[
 a_3=1+xy,\qquad a_5=1+xz,\qquad a_7=1+yz.
\]
Here $n=24$, and the correct reflection dimension is
\begin{equation}\label{ns:eq:counter-dimension}
 24\cdot8+\frac{4+14+14+4}{2}=210.
\end{equation}
The naive colength substitution would instead give $208$. Both integral
Tate parity sums equal $18$, hence
\begin{equation}\label{ns:eq:counter-smith}
 C_\pm\simeq\mathbb Z^{192}\oplus(\mathbb Z/2\mathbb Z)^{18}
\end{equation}
as abelian groups. The failed extrapolation would predict only sixteen
torsion factors. This is a counterexample to a possible extension, not an
error in the present manuscript: its separability restriction explicitly
excludes the example.

\subsection{Independent checks and the remaining module problem}

The accompanying program \texttt{verify\_nonseparable\_jets.py} performs exact
binary elimination by two independent constructions. One builds the Koszul
matrices from wedge bases. The other expands the original prime rows of
\eqref{ns:eq:distribution} and the reflection rows in point/monomial
coordinates; it imports neither a normal form nor a Koszul calculation.
The resulting raw dimensions are
\begin{center}
\begin{tabular}{lrrr}
\hline
Example & Generators & Relation rank & Quotient dimension\\
\hline
Mixed parameters, $q=15$ &135&91&44\\
Three-generator obstruction, $q=105$ &840&630&210\\
Redundant third parameter, $q=105$ &945&713&232\\
\hline
\end{tabular}
\end{center}
Their unreflected dimensions are $72,384,432$, as required by the all-ring
basis theorem. Omitting reflection rows supplies a corruption control:
every reported dimension changes. These checks verify the concrete examples;
the preceding proofs establish the theorem for all permitted rings and
parameters.

For three active generators, Frobenius duality always reduces the homology
dimension profile to $(\ell,h_1,h_1,\ell)$, so the reflection defect is
$\ell+h_1$. Computing the syzygy invariant $h_1$ is the first unresolved step
beyond the two-generator law. Further useful questions are to classify ideals
for which the binomial colength profile still holds, determine when the
positive and negative torsion modules are isomorphic, and characterize the
deformations realizable by actual arithmetic spectral jets. Those questions
must retain the distinction between module structure, total dimension, and
relations among analytic values.
